% Frontiers Brief Research Report; official Harvard class, one-column review layout.
% Compile with pdfLaTeX + BibTeX. Main document: main.tex.
\documentclass[utf8]{FrontiersinHarvard}
\usepackage{url,hyperref,lineno,microtype,booktabs}
\usepackage[singlespacing]{setspace}
\usepackage{accsupp}
\hypersetup{hidelinks}
% Allow the official running head its measured height without moving the text block.
\setlength{\headheight}{20pt}
\setlength{\headsep}{15pt}
\linenumbers
\def\keyFont{\fontsize{8}{11}\helveticabold}
% Author details restored from the supplied Ver 2.1 manuscript.
\def\firstAuthorLast{Hossain {et~al.}}
\def\Authors{Tanzir Hossain\,$^{1}$, Rajib Rana\,$^{2,*}$, Prabal Datta Barua\,$^{2,7}$, Abu Ali Ibn Sina\,$^{3}$, Niall Higgins\,$^{2,4}$, Pascal Elahi\,$^{8}$, Robert Sang\,$^{2}$, and Bj\"orn W. Schuller\,$^{5,6}$}
\def\Address{$^{1}$BRAC University, Dhaka, Bangladesh\\
$^{2}$University of Southern Queensland, Springfield Central, QLD, Australia\\
$^{3}$School of Biotechnology and Biomolecular Sciences, UNSW Sydney, Sydney, NSW, Australia\\
$^{4}$Mental Health and Specialist Services, West Moreton Health, Brisbane, QLD, Australia\\
$^{5}$GLAM--Group on Language, Audio, and Music, Imperial College London, London, United Kingdom\\
$^{6}$Chair of Health Informatics, Technical University of Munich (TUM), Munich, Germany\\
$^{7}$Cogninet AI, Brisbane, QLD, Australia\\
$^{8}$Pawsey Supercomputing Research Centre, Kensington, WA, Australia}
\def\corrAuthor{Rajib Rana}
\def\corrEmail{Rajib.Rana@unisq.edu.au}
\newcommand{\AuthorContributions}{TH contributed to methodology, software development, formal analysis, data curation, visualization, and preparation of the original draft. RR contributed to conceptualization, methodology, project administration, supervision, funding acquisition, interpretation of results, and writing and editing of the manuscript. PDB contributed to methodology, validation, interpretation of results, and review and editing of the manuscript. AAIS contributed to interpretation of the drug-response model and biological context, and review and editing of the manuscript. NH contributed to interpretation of the study, critical review, and editing of the manuscript. PE contributed to high-performance and quantum computing methodology, computational resource support, technical validation, and review and editing of the manuscript. RS contributed to supervision, research direction, interpretation of the findings, and review and editing of the manuscript. BWS contributed to conceptualization, methodological guidance, supervision, interpretation of the quantum-computing results, and review and editing of the manuscript.}
\newcommand{\FundingStatement}{This work was supported by computational resources provided by the Pawsey Supercomputing Research Centre through the Setonix-Q Pilot Scheme under project NCMAS-2026-66, “Closing the Causality Gap: A Computational Framework for Validating Allosteric Drug Mechanisms.” The allocation provided access to Pawsey’s Setonix high-performance computing infrastructure and quantum-enabled computing resources, including access to quantum processing resources through the Quantum Hub Portal and AWS Braket.}
\newcommand{\ConflictStatement}{The authors declare that the research was conducted in the absence of any commercial or financial relationships that could be construed as a potential conflict of interest.}
\newcommand{\DataStatement}{The drug-response data used in this study were obtained from the publicly available AstraZeneca--Sanger Drug Combination DREAM Challenge dataset. The processed optimization instances, quantum hardware output data, noiseless simulation results, classical baseline results, parameter-sensitivity analyses, and scripts used to reproduce the analyses will be made publicly available in a research repository upon publication. Quantum hardware task identifiers and sufficient metadata to reproduce the reported analyses will also be provided, subject to the access conditions of the relevant quantum-computing platforms.}

\newcommand{\TableOneRows}{%
6 & 3 & 2 & 30 & 20 & 176/200 & 88.0 (82.8--91.8) & 5/200  & 100 \\
8 & 4 & 2 & 56 & 56 & 146/200 & 73.0 (66.5--78.7) & 5/200  & 100 \\
9 & 3 & 3 & 72 & 54 & 126/200 & 63.0 (56.1--69.4) & 0/200  & 100 \\
10 & 5 & 2 & 90 & 120 & 141/200 & 70.5 (63.8--76.4) & 0/200  & 100 \\
12 & 4 & 3 & 132 & 162 & 59/200 & 29.5 (23.6--36.2) & 0/200  & 100 \\
15 & 5 & 3 & 210 & 360 & 130/200 & 65.0 (58.2--71.3) & 0/200  & 98 \\
}

\newcommand{\RecoveredRows}{%
6 & 3 & 2 & 07-03 03:33 & 500 & 88.4 (85.3--90.9) & 6\\
6 & 3 & 2 & 07-11 07:10 & 200 & 85.5 (80.0--89.7) & 5\\
8 & 4 & 2 & 07-11 07:12 & 200 & 72.0 (65.4--77.8) & 5\\
10 & 5 & 2 & 07-11 07:19 & 200 & 69.5 (62.8--75.5) & 4\\
12 & 6 & 2 & 07-11 07:25 & 200 & 72.5 (65.9--78.2) & 4\\
14 & 7 & 2 & 07-11 07:38 & 200 & 62.0 (55.1--68.4) & 3\\
15 & 5 & 3 & 07-11 07:48 & 200 & 54.5 (47.6--61.3) & 0\\
18 & 6 & 3 & 07-11 06:06 & 200 & 42.0 (35.4--48.9) & 0\\
18 & 6 & 3 & 07-11 07:55 & 200 & 49.5 (42.6--56.4) & 1\\
18 & 6 & 3 & 07-12 13:37 & 200 & 46.0 (39.2--52.9) & 1\\
20 & 4 & 5 & 07-11 08:09 & 200 & 20.0 (15.0--26.1) & 0\\
20 & 4 & 5 & 07-12 13:47 & 200 & 27.0 (21.3--33.5) & 0\\
21 & 7 & 3 & 07-11 06:09 & 200 & 44.0 (37.3--50.9) & 1\\
21 & 7 & 3 & 07-11 08:02 & 200 & 40.5 (33.9--47.4) & 0\\
21 & 7 & 3 & 07-12 13:40 & 200 & 35.0 (28.7--41.8) & 0\\
25 & 5 & 5 & 07-11 08:17 & 200 & 13.0 (9.0--18.4) & 0\\
25 & 5 & 5 & 07-12 13:54 & 200 & 15.0 (10.7--20.6) & 0\\
30 & 6 & 5 & 07-11 08:23 & 200 & 6.5 (3.8--10.8) & 0\\
35 & 7 & 5 & 07-11 08:33 & 200 & 3.5 (1.7--7.0) & 0\\
}

\newcommand{\Fset}{\mathcal{F}}
\newcommand{\altfigure}[3]{\BeginAccSupp{method=escape,ActualText={#2}}\includegraphics[width=#3]{#1}\EndAccSupp{}}


\firstpage{1}
\title[Warm-Start QAOA Benchmarking]{Feasibility and optimum recovery in warm-start quantum optimization for a drug-response model on a trapped-ion processor}



\begin{document}

\onecolumn
% \firstpage{1}

% \title[Warm-Start QAOA Benchmarking]
% {Benchmarking warm-start QAOA for constrained optimization on trapped-ion hardware}

\author[\firstAuthorLast]{\Authors}

\address{}
\correspondance{}
\extraAuth{}

{\setlength{\parindent}{0pt}\maketitle}
% TeXcount body text including author-date citations.
\newcommand{\MainTextWords}{5042}

\noindent\textbf{Article type:} Brief Research Report\quad
\textbf{Main-text words:} \MainTextWords\quad
\textbf{Figures:} 2\quad\textbf{Tables:} 2

\begin{abstract}
\section{}
We evaluate whether warm-start quantum approximate optimization (QAOA) can find the best drug-and-dose assignments in a model built from measured individual and pairwise drug responses. A warm start uses classical information to guide the circuit's initial state. On the IonQ Forte-1 processor, this approach returned valid assignments more often than uniformly random binary strings. However, producing valid assignments did not consistently lead to finding the best assignment. Across twelve reference circuits, ideal, noiseless QAOA improved the probability of finding an optimum over equally likely sampling of valid assignments in only four cases. Hardware often performed below the noiseless prediction for the same circuit, while classical greedy search solved every hardware-associated model within 200 score evaluations. We analyzed 25 completed hardware tasks containing 5,300 measurements from 18 circuits and 14 problem instances. These used 6--35 qubits and at most 4,900 valid assignments, allowing us to check every valid assignment and determine the exact optima. Noiseless calculations cover the original six circuits and six wider circuits. At 35 qubits, the ideal probability of a valid output was 35.85\%, whereas hardware returned only 7 valid outputs in 200 measurements; the ideal optimum probability was also below both sampling controls. Further simulations improved optimum probability over uniform valid sampling in 28 of 35 CAMA-1 compound panels and none of four 647-V panels. Classical simulated annealing solved all these panels for every tested random seed. A Grover mixer, which keeps the circuit within valid assignments, preserved validity and improved optimum probability over uniform valid sampling in all ten tested models.


\tiny
\keyFont{\section{Keywords:} quantum approximate optimization, warm start, trapped-ion quantum computing, feasibility, classical baselines, drug-response modeling, constrained optimization}
\end{abstract}

\section{Introduction}
Selecting a combination of anticancer compounds requires accounting for the response to each compound and for interactions between them. Dose choices add another set of decisions. Experimental combination screens provide measurements for constructing such models, but a computational score still needs to be distinguished from a validated treatment response \citep{menden2019}. Here, drug-and-dose selection provides a measured-data setting in which to examine how reliably a quantum optimization method finds high-scoring, valid choices.

A constrained optimization method must satisfy two goals: return a choice that obeys the rules and find a choice with a high score. We call these properties feasibility and solution quality. A quantum circuit can improve feasibility over random binary strings while doing little to improve solution quality. We examine this distinction using archived executions of warm-start QAOA on a trapped-ion processor. The task is to select compounds and doses in a drug-response model. We compare the circuit's starting distribution, its ideal behavior without noise, and its hardware outputs with the exact best assignments and with classical search.

QAOA alternates two operations: cost evolution, which applies score-dependent phases to quantum states, and mixing, which changes their amplitudes and therefore their probabilities when measured \citep{farhi2014}. Warm-start methods use classical information to choose the initial state and, in some approaches, the mixer \citep{egger2021}. They allow the quantum search to begin with information about the problem already encoded in its sampling probabilities. Constraint-preserving mixers address a different concern: they keep the circuit within the set of valid assignments \citep{hadfield2019}. These approaches motivate separate checks of what the starting state contributes, what the quantum circuit adds, and whether its outputs obey the constraints.

Earlier studies reported warm-start improvements for portfolio selection and MaxCut, a graph-partitioning problem \citep{egger2021}, and trapped-ion optimization using up to 40 qubits \citep{pagano2020}. In constrained text summarization, \citet{niroula2022} compared feasibility and solution quality with uniform feasible sampling at 14 and 20 qubits, showing why a high rate of valid outputs alone is an incomplete performance measure. \citet{ramos2026} studied drug-combination models expressed as quadratic unconstrained binary optimization problems, or QUBOs, using simulated quantum and classical annealing. Together, these studies motivate controls that distinguish valid sampling from successful optimization and include classical methods that actively search for high scores \citep{koch2026,sankar2024}. Our contribution is to connect those controls with archived drug-dose circuits, exact feasible optima, and noiseless predictions for the same circuits, so that improvements due to initialization can be separated from circuit behavior and hardware execution.

Our model uses data from the AstraZeneca--Sanger Drug Combination DREAM study \citep{menden2019}. We compare 25 hardware executions with the initial sampling distributions, noiseless calculations of the same circuits, exact optima, and classical search. Ideal QAOA improves optimum probability over its initial distribution in eleven of twelve reference circuits, but exceeds uniform sampling of valid assignments in only four. Four additional compound panels test whether the gains extend to problems with different optima. Further simulations cover all CAMA-1 panels, a second cell line, other score weights, and a mixer that preserves feasibility. Greedy search solves every hardware-associated problem at the archived score weight. Even the largest instance, encoded with 35 qubits, has only 4,900 valid assignments and can be checked exhaustively. We therefore report both circuit width, which measures qubit resources, and the number of valid assignments, which measures the choices the optimization must search.

\section{Method}
\subsection{Source data and instance construction}
We use the AstraZeneca--Sanger Drug Combination DREAM dataset to construct optimization problems whose coefficients come from experimental drug responses \citep{menden2019}. In the original screen, researchers exposed laboratory-grown cancer cells to individual anticancer compounds and pairs of compounds at different concentrations, then measured cell viability. The dataset's single-agent dose-response parameters describe how viability changes with concentration. Its pairwise synergy scores describe interactions between compounds. We use these two types of measurement to construct the individual-compound and pair-interaction contributions to the model score.

The archived model contains seven compounds tested on the CAMA-1 cancer cell line, a population of cancer cells maintained for laboratory experiments. Every one of the 21 possible pairs has a recorded synergy observation, so we do not need to estimate missing pair coefficients. The source pipeline retains observations that pass the dataset's quality check, averages repeated synergy scores for each pair, and takes the median single-agent dose-response parameters for each compound. These summaries provide one set of coefficients for each compound and pair; we retain the archived aggregation rules to reproduce the models used on hardware.

The dataset identifies the compounds by their molecular targets, the proteins on which they act. AKT\_1, MAP2K\_1 and MTOR\_1 refer to compounds targeting AKT, MEK1/2 and mTOR, respectively; these are enzymes involved in cell signaling. PIK3CB\_PIK3CD refers to the beta and delta forms of phosphoinositide 3-kinase, another signaling enzyme. BCL2\_BCL2L1 refers to proteins that regulate programmed cell death, and FASN refers to fatty acid synthase, an enzyme that makes fatty acids. SLC16A4 is the dataset's compound label for the solute-carrier targets SLC16A3 and SLC16A4, which belong to a membrane-transporter family \citep{ncbigene2026}. We retain the source identifiers so readers can trace the model coefficients to the experimental records.

The hardware-associated problems offer three to seven candidate compounds and two, three or five nonzero dose choices per compound. As described below, an allowed assignment selects two or three of those compounds and one dose for each. We keep the problems small enough to evaluate every allowed assignment and determine the exact best model score. That gives us a reference against which to assess each optimization method. More candidate compounds and dose choices increase both the number of assignments to check and the qubit resources required by this encoding.

\subsection{Dose-assignment representation and score}
For $n$ candidate compounds and $D$ nonzero dose choices per compound, we use $q=nD$ binary variables. Compound indices $i,j$ run from 1 to $n$, and dose indices $d,e$ run from 1 to $D$. The variable $x_{id}\in\{0,1\}$ is one when compound $i$ is selected at dose $d$ and zero otherwise. We write $x$ for the complete assignment of these variables. Each compound has a block of $D$ variables and all zeros means that the compound is absent, and a single one selects its dose. Two or more ones in a block would assign multiple doses to the same compound and are invalid. A feasible assignment selects two or three compounds, at most one dose for each, and at least two recorded classes. These classes group compounds by their mechanism labels. Every compound in each included instance has a distinct class label, so selecting at least two compounds already satisfies the class requirement. These rules govern the structure of an assignment; the archived model includes no exposure cap or clinical safety constraint.

We maximize the scalar objective $f(x)$, the model score of assignment $x$. The coefficient $h_{id}$ describes the individual effect of compound $i$ at dose $d$, $t_{id}$ is its dose-based penalty proxy, and $\lambda$ weights that proxy. The coefficient $J_{id,je}$ describes the interaction between compound $i$ at dose $d$ and compound $j$ at dose $e$. The score is
\begin{equation}
 f(x)=\sum_{i,d}(h_{id}-\lambda t_{id})x_{id}
       +\sum_{i<j}\sum_{d,e}J_{id,je}x_{id}x_{je},\qquad\lambda=0.6.
 \label{eq:score}
\end{equation}
The first sum covers all compound-dose choices. The restriction $i<j$ in the second sum counts each distinct compound pair once, while $d$ and $e$ cover its dose choices. The coefficient $h_{id}$ estimates inhibition by compound $i$ at dose $d$ using a fitted Hill dose-response curve, which relates concentration to cell viability. Dose choices are fractions of each compound's maximum tested concentration. The original series uses $(0.1,1)$ or $(0.1,0.32,1)$; the larger instances also use the archived five-level grids. The interaction term $J_{id,je}$ scales the measured synergy coefficient for compounds $i$ and $j$ by the product of their dose fractions. This scaling assumes how synergy changes with dose; it is not fitted to a measured dose-response surface. For a three-compound selection, the score includes three individual effects and the three pair interactions. It therefore estimates the three-compound score from individual and pairwise measurements.

The term $t_{id}$ is a dose-based proxy weighted by mechanism class, and $\lambda=0.6$ sets its contribution to the score. It is not a measured toxicity value. We retain this weight and the archived coefficients to analyze the models used in the hardware runs, and test other weights separately. The value 0.6 is an archived modeling choice, not a clinically calibrated trade-off.

For larger instances, we construct the two- and three-compound assignments directly and separately decode every observed binary string. This avoids storing all $2^q$ possible binary states. Since every compound has a distinct mechanism class, the number of feasible assignments is
\begin{equation}
 |\Fset|=\binom{n}{2}D^2+\binom{n}{3}D^3.
 \label{eq:count}
\end{equation}
Here, $|\Fset|$ is the number of feasible assignments, and $\binom{n}{k}$ counts the ways to choose $k$ distinct compounds from $n$. The first term counts all two-compound selections and their dose choices; the second counts all three-compound selections and their dose choices. The original six instances contain 20--360 feasible assignments. The expanded set reaches 4,900 at 35 qubits. We find each exact optimum by evaluating Equation~\ref{eq:score} for every feasible assignment.

\subsection{Quantum hardware and execution}
We analyze circuits executed on the IonQ Forte-1 trapped-ion quantum processor through Amazon Braket. Forte stores qubits in trapped ions, controls them with lasers, and supports interactions between any pair of qubits \citep{chen2024}. The cost Hamiltonian is the operator that encodes the score and constraint penalties. Its two-qubit terms require operations between the corresponding qubit pairs. Forte's all-to-all connectivity supports these operations without restricting the model to a nearest-neighbor coupling pattern; it does not mean that all pair interactions execute simultaneously. Our experiments use 6--35 qubits to represent problems with 20--4,900 feasible assignments, with one qubit for each compound-dose choice. The submitted programs use single-qubit rotations and controlled-NOT, or CNOT, gates. They contain 30--1,190 CNOT gates. These are counts in the submitted circuits; counts after compilation into the processor's native operations are unavailable.

A shot is one execution and measurement of a circuit. Each shot returns a binary string that we decode into compound and dose choices. We analyze 25 hardware executions of 18 distinct circuits representing 14 optimization instances. Twenty-four executions contain 200 shots and one contains 500, giving 5,300 measurements in total. We keep repeated executions of identical circuits separate, with six 6--15-qubit circuits forming the original reference series (Table~\ref{tab:hardware}); the additional instances reach 35 qubits and 4,900 feasible assignments (Table~\ref{tab:expanded}). We compare each execution with classical controls and noiseless calculations of the same submitted circuit. Circuit details and task identifiers are provided in the supplement.

\subsection{Executed circuit and noiseless reconstruction}
The archived circuit begins in a product state: each qubit is prepared independently, so the probability of an initial binary string is the product of its individual bit probabilities. We index qubits by $v,w\in\{0,\ldots,q-1\}$, using $v=(i-1)D+(d-1)$ for compound-dose variable $x_{id}$; $x_v$ is the same bit in this flattened ordering. The symbol $|0\rangle$ denotes a qubit's zero state, and $R_y(\theta_v)$ rotates qubit $v$ about its $y$ axis through angle $\theta_v$. The tensor product $\bigotimes$ combines the independently prepared qubit states. The initial state $|\psi_0\rangle$ and the probability $P_0(x)$ of measuring assignment $x$ are
\begin{equation}
 |\psi_0\rangle=\bigotimes_{v=0}^{q-1}R_y(\theta_v)|0\rangle,
 \qquad P_0(x)=\prod_v c_v^{x_v}(1-c_v)^{1-x_v},\quad
 c_v=\sin^2(\theta_v/2).
 \label{eq:prior}
\end{equation}
The product $\prod_v$ runs over all qubits. The rotation angle $\theta_v$ sets the probability $c_v$ that bit $v$ is one, so $1-c_v$ is its probability of being zero. Here, $c_v$ denotes a bit probability, not a physical dose. We refer to this initial distribution as the prior. The archived circuit contains one QAOA layer. It first applies the diagonal Ising cost operation, which assigns phases according to the encoded cost, then applies a separate warm-start mixer to each qubit:
\begin{equation}
 |\psi\rangle=\left[\bigotimes_v R_y(\theta_v)R_z(2\beta)R_y(-\theta_v)\right]
 \exp\!\left[-i\gamma\left(\sum_v a_vZ_v+\sum_{v<w}b_{vw}Z_vZ_w\right)\right]|\psi_0\rangle.
 \label{eq:circuit}
\end{equation}
In Equation~\ref{eq:circuit}, $|\psi\rangle$ is the state after the QAOA layer. The Pauli-$Z$ operator $Z_v$ has eigenvalues +1 and -1 for bit values zero and one. The real coefficients $a_v$ and $b_{vw}$ weight the single-qubit and pairwise cost terms, including constraint penalties; $v<w$ counts each qubit pair once. The angles $\gamma$ and $\beta$ control the cost and mixing operations, respectively, while the angles $\theta_v$ remain fixed by initialization. $R_z$ is a rotation about the $z$ axis, $\exp$ is the operator exponential, and the $i$ in its exponent is the imaginary unit, satisfying $i^2=-1$. Operator products act from right to left. All rotation angles are expressed in radians. We recover the coefficients and angles from the archived execution scripts and check them against the submitted programs. The cost includes penalties for invalid assignments. Mixing uses single-qubit rotations, while the CNOT--$R_z$--CNOT blocks implement the pairwise cost interactions. These circuits do not use an XY exchange mixer, which exchanges excitations between qubits, and they do not guarantee that an output is feasible.

We retain the parameters used in the archived hardware executions. For the original six circuits, we compare a NumPy gate-by-gate simulator and an independent calculation of Equation~\ref{eq:circuit} with PennyLane. We extend the latter calculation to four circuits using 18--25 qubits. To reduce memory demands, we calculate phases in chunks and update pairs of amplitudes in place. At each of these widths, the repeated hardware tasks use the same circuit.

We also calculate the full statevector, the list of quantum amplitudes for all binary states, at 30 qubits and 35 qubits on a Pawsey node with sufficient memory. These circuits encode problems with 2,875 and 4,900 feasible assignments, respectively. The calculation uses single-precision complex amplitudes, initializes them in Gray-code order, in which successive binary states differ by one bit, and applies local mixers in parallel. At 25 qubits, the feasible probability differs from a double-precision calculation by less than $10^{-7}$. The accelerated implementation also reproduces the distributions of all six original circuits. We preserve the recorded qubit order, with the submitted program's first qubit, q[0], corresponding to the first outcome bit. The supplement gives precision checks, memory requirements, and details of transferred parameters.

\subsection{Classical controls and performance measures}
We compare the quantum results with several classical controls. Uniform sampling over all $2^q$ binary strings gives every encoded state the same probability, including invalid states. Its probability of feasibility is exactly $|\Fset|/2^q$. The warm-start-only control samples each bit independently using Equation~\ref{eq:prior}. This can be done classically and measures what the starting distribution achieves before the cost and mixing operations.

Uniform feasible sampling gives every valid assignment the same probability without evaluating its score. To sample this way, we first choose a combination size $m\in\{2,3\}$ with probability proportional to $\binom{n}{m}D^m$, then choose a compound subset and its doses uniformly. Weighting the size choice by the number of assignments at that size makes all feasible assignments equally likely. Every output is valid. We also use exhaustive enumeration as an exact solver, evaluating the score once for every feasible assignment.

For each distribution, we calculate the probability $P_{\Fset}$ of a valid assignment and the probability $P_\star$ of any assignment with the globally best feasible score. When these probabilities are known analytically, the probability of obtaining at least one optimum in $B$ independent draws is
\begin{equation}
 P_{\mathrm{hit}}(B)=1-(1-P_\star)^B,\qquad B=200.
 \label{eq:budget}
\end{equation}
Here, $P_{\mathrm{hit}}(B)$ is the probability of at least one optimum in $B$ draws; $P_\star$ is the optimum probability on a single draw. We use $B=200$ to match the usual hardware shot count. These probabilities are calculated directly from the distributions, so the analytic sampling-control results do not depend on a Monte Carlo seed.

We check the recorded hardware probabilities by multiplying each one by its task's shot count. Every result is within $10^{-7}$ of an integer, and the counts sum to the declared 200 or 500 shots. We calculate feasibility using all shots as the denominator, including invalid outputs. For hardware proportions, we report pointwise 95\% Wilson confidence intervals to express the sampling uncertainty in each estimate \citep{wilson1927}.

To compare the quality of valid outputs, we rescale scores to $u(x)=[f(x)-f_{\min}]/[f_\star-f_{\min}]$, where $f_{\min}$ and $f_\star$ are the worst and best scores over $\Fset$. Thus, zero is the worst feasible score and one is the optimum. We report the mean $u$ among valid hardware shots and the best-sample gap $1-\max u$, where the maximum is over valid outputs in the observed batch. This gap measures how far the best observed score falls below the optimum as a fraction of the feasible score range. Invalid outputs are counted separately through feasibility. These summaries describe the observed batch.

Our greedy search starts with an empty assignment and adds the best available compound-dose choice until two compounds have been selected. It then considers dose changes, additions, deletions, and compound swaps that preserve feasibility, choosing the change that improves the score most. Previously calculated scores are cached so that they are not evaluated again. The search stops after 200 unique score evaluations, including partial assignments during construction, or when no neighboring assignment improves the score. The fixed cap provides a common score-evaluation budget for the two classical searches, without implying equal runtime or equal cost to 200 quantum shots. We compare its result with the exact optimum only after the search has finished.

We also run simulated annealing, which can accept a temporary decrease in score to explore beyond a local optimum. It starts from a random feasible assignment, uses the same permitted changes as greedy search, and gradually reduces its acceptance of worse choices through geometric cooling. We use seeds 0--99, with a limit of 200 unique score evaluations per run. These repeated starts measure dependence on random initialization and moves; the choice of 100 runs is a fixed experimental budget. We report the fraction of runs that reach the exact optimum and, among successful runs, the median number of evaluations needed to first reach it. The supplement specifies the proposal limit and temperature schedule.

We retain the archived score weight $\lambda=0.6$ for the main comparisons and test weights from 0.3 to 0.9 in steps of 0.15. At each weight, we recompute the exact optima, rescore the fixed sampling distributions, and rerun both classical searches. The hardware-associated models keep their optima and the rankings of the sampling controls across this range, while rankings change when QAOA is retrained for the additional panels (Supplementary Table~S9).

Finally, we simulate one-layer QAOA with a Grover mixer for the six original instances and four original panels. The circuit starts with equal amplitudes on all feasible assignments. It uses the score without constraint penalties and a mixer that applies a selective phase shift to that initial feasible state \citep{bartschi2020}. This construction keeps feasibility at one and tests whether avoiding invalid assignments improves optimum recovery. We train both circuit angles using the same three random seeds and Adam settings as in the panel simulations described next.

\subsection{Panel selection and initialization}
For the additional simulations, we select four distinct four-compound panels from the 35 possible panels in the seven-compound CAMA-1 model. Panel size, dose count, circuit width, feasible-set size, score weight, and constraints are kept fixed so that only compound membership changes. The four panels are sampled uniformly without replacement using a fixed NumPy seed \citep{numpy2026} and are frozen before optimization. For each panel, we construct the penalized cost Hamiltonian and warm-start product distribution, using a fixed one-layer circuit with local warm-start mixers. Adam optimizes only the two shared QAOA angles. Each panel is trained with seeds 0, 1, and 2 for 60 updates at a learning rate of 0.1, and we retain the final parameters from every run rather than selecting runs by optimum probability.

We then extend the same procedure to all 35 CAMA-1 panels, four 647-V panels selected by measurement coverage, and the original four panels retrained across different score weights. The supplement reports panel membership, selection rules, and complete results. For reproducibility, the repository records the code, coefficients, frozen panel membership, seeds, parameter values, training histories, software environments, and source-file hashes. Because small numerical differences across platforms can alter optimization trajectories, reproducibility is assessed numerically rather than assumed to be bit-for-bit identical. Saved hardware outcomes allow the reported analysis to be repeated without rerunning the quantum processor.

\section{Results}
\subsection{Original six-task reference: feasibility}
Hardware produces valid assignments more often than uniformly random binary strings in every original reference instance. Hardware feasibility ranges from 29.5\% to 88.0\%, compared with 1.10--31.25\% for random binary strings. The submitted circuits contain 30--210 CNOT gates (Table~\ref{tab:hardware}).

The largest relative gain arises where very few binary strings represent valid assignments. At 15 qubits, only 1.10\% of binary strings are feasible. Hardware improves on that probability by a factor of 59.2, yet more than a third of its outputs are still invalid. Uniform feasible sampling produces a valid assignment on every draw (Figure~\ref{fig:feasibility}A).

Hardware also exceeds the warm-start-only control in five of the six cases. The exception is the 12-qubit circuit, where the ideal circuit already lowers feasibility from the prior's 32.51\% to 31.18\%. Across the reference circuits, ideal feasibility ranges from 31.18\% to 89.96\%, while the prior ranges from 31.68\% to 62.42\% (Supplementary Table~S1).

\subsection{Optimum recovery in the reference subset}
Hardware finds an optimum in only two of the six reference batches. Independent enumeration identifies the same unique optimum in all six models: the same three compounds at their highest encoded dose, with a model score of approximately $1.769$. We round the value in the prose for readability; the archived numerical results retain fuller precision for computational checks. These digits describe the model calculation and do not imply equivalent precision in measured biological effects. Hardware returns this assignment five times at 6 qubits and five times at 8 qubits, and never in the other four batches. For a batch with no observed optimum, the Wilson interval still has an upper endpoint of approximately 1.88\% for the per-shot probability. Observing zero successes therefore does not establish that success is impossible (Table~\ref{tab:hardware}).

The 15-qubit batch shows how feasibility and optimum recovery can differ. It contains 130 valid outputs and no optimum. The ideal circuit's per-shot optimum probability is 0.247\%, compared with 0.144\% for the prior and 0.278\% for uniform feasible sampling. Over 200 independent draws, these probabilities give calculated success rates of 39.0\%, 25.0\%, and 42.7\%, respectively. Uniform sampling of all binary strings gives a success rate of only 0.61\% (Figure~\ref{fig:optimum}).

Ideal QAOA improves optimum probability over initialization in every original reference circuit, but exceeds uniform feasible sampling only at 8 and 10 qubits. At 10 qubits, the ideal circuit predicts a 94.0\% chance of finding an optimum within 200 draws, while hardware finds none. A later hardware execution uses different parameters, so it cannot resolve the discrepancy for the original circuit (Figure~\ref{fig:optimum}; Supplementary Table~S2).

\subsection{Solution quality and a constructive search baseline}
Greedy local search finds the exact optimum for all six reference instances within 21--81 unique score evaluations. For hardware, the gap between the best observed score and the optimum ranges from zero to 39.5\% of the feasible score range. The four batches that miss the optimum have gaps of 18.8--39.5\% (Supplementary Table~S4 and Figure~S1).

Across all 14 hardware-associated models at $\lambda=0.6$, simulated annealing reaches the optimum in 15--100\% of runs within the evaluation limit. For the original six instances, success is 98--100\%, with median first-hit counts of 11--87.5 evaluations among successful runs. Greedy search solves all 14 models in 21--200 evaluations, while exact enumeration requires 20--4,900. In the 35-qubit instance, which has 4,900 feasible assignments, annealing succeeds for 15 of 100 seeds and greedy search first evaluates an optimal assignment after 153 unique score evaluations and completes its 200-evaluation budget (Table~\ref{tab:hardware}; Supplementary Table~S8).

\subsection{Expanded hardware records and repeated executions}
Hardware does not consistently recover an optimum in the additional instances. It succeeds in the 10--14-qubit cases, in two of three 18-qubit executions, and in one of three 21-qubit executions. It finds no optimum in any instance with five dose levels, including the 35-qubit encoding with 4,900 feasible assignments (Table~\ref{tab:expanded}).

Some repeated executions of the same circuit differ appreciably. Feasibility ranges from 35.0\% to 44.0\% at 21 qubits and from 42.0\% to 49.5\% at 18 qubits. The two identical-circuit 6-qubit runs are closer, at 88.0\% and 88.4\%, with 200 and 500 shots, respectively. We keep circuits with changed parameters or compound membership as separate records (Supplementary Figure~S2).

At 18, 21 and 25 qubits, hardware feasibility is below the prediction for the same noiseless circuit. Every pointwise hardware confidence interval at these settings excludes the ideal value. Both 20-qubit intervals include it. Ideal feasibility is 73.16\% at 18 qubits and 73.44\% at 21 qubits. The ideal 20- and 25-qubit circuits already reduce feasibility below their starting distributions; hardware introduces a further shortfall at 25 qubits (Figure~\ref{fig:feasibility}B; Supplementary Tables~S6--S7).

At all four of these settings, ideal QAOA improves optimum probability over initialization, but exceeds uniform feasible sampling only at 18 and 21 qubits. Ideal optimum probability ranges from 0.0239\% to 0.1666\%. Among valid outputs, mean normalized quality is 39.0--51.0\% ideally and 37.7--50.0\% on hardware. These means exclude invalid outputs (Supplementary Tables~S6--S7).

The two widest circuits combine weak ideal optimization with a further loss of feasibility on hardware. At 30 qubits, with 2,875 feasible assignments, ideal feasibility is 33.39\%, close to the prior's 33.64\%. Hardware returns only 13 valid outputs in 200 shots. At 35 qubits, with 4,900 feasible assignments, ideal feasibility is 35.85\%, close to the prior's 35.75\%, while hardware returns 7 valid outputs in 200 shots. The latter batch has a Wilson interval of 1.7--7.0\% for feasibility and a best-sample normalized gap of 40.4\%. Both ideal optimum probabilities fall below uniform feasible sampling. At 35 qubits, the ideal optimum probability of 0.01197\% also falls below the prior's 0.01217\% (Figure~\ref{fig:feasibility}B; Supplementary Tables~S6--S7).

\subsection{Results for the additional panels}
Ideal QAOA improves optimum probability over both sampling controls in every run of the original four-panel simulation. Each panel has 162 feasible assignments and a different unique optimum. Across the 12 panel-and-seed combinations, ideal optimum probability is 2.23--3.43\%, compared with 1.26--1.48\% for initialization and 0.617\% for uniform feasible sampling. Greedy search solves each panel in 36--50 evaluations (Supplementary Table~S3).

The expanded panel study shows that this gain depends on the model. Ideal QAOA exceeds uniform feasible sampling for all three seeds in 28 of 35 CAMA-1 panels, or 80\%, and in none of the four 647-V panels. Greedy search solves all CAMA-1 panels and three of four 647-V panels, or 75\%. Simulated annealing solves every panel for all 100 seeds (Supplementary Table~S11 and Figure~S3).

Changing the score weight leaves the archived sampling comparisons unchanged but affects retrained QAOA and classical search. All 14 hardware-associated models retain their respective optima across the weight sweep. Greedy search recovers the optimum in 61 of 70 model-and-weight combinations. Across the original four panels and five weights, retrained QAOA exceeds uniform feasible sampling in 52 of 60 runs and its initial product distribution in 59 of 60 (Supplementary Table~S9).

\subsection{Feasibility-preserving mixer results}
The Grover mixer keeps every output feasible and improves optimum probability over uniform feasible sampling in every tested run. Across the six original instances and four original panels, optimum probability ranges from 1.17\% to 27.24\%. It exceeds the archived penalty-QAOA probability on all six reference instances, but falls below the trained penalty-QAOA probability on all four panels. Here, penalty-QAOA refers to the circuits that encode constraints through cost penalties. Greedy search solves all ten models (Supplementary Table~S10).

\section{Discussion}
Warm-start QAOA improves some sampling probabilities in this model, but a higher rate of valid outputs alone does not establish effective optimization. Ideal QAOA improves optimum probability over initialization in eleven of twelve reference circuits and over uniform feasible sampling in only four. Across the wider panel study, the gain over uniform feasible sampling occurs in 28 of 35 CAMA-1 panels and none of four 647-V panels. At the archived weight, greedy search solves all hardware-associated models and all CAMA-1 panels. Simulated annealing also solves the 647-V panel that greedy search misses. These results do not support a claim of computational advantage.

Noiseless calculations help separate limitations of the circuit design from losses during hardware execution. The ideal 20- and 25-qubit circuits slightly reduce feasibility relative to initialization. At 18 and 21 qubits, hardware loses much of the feasible probability predicted by the ideal circuit; at 25 qubits, it adds to an existing circuit limitation. The 30- and 35-qubit references, with 2,875 and 4,900 feasible assignments, show little ideal optimization gain and a substantial further loss of feasible outputs on hardware. Identifying the cause of an execution discrepancy would require state-preparation controls, calibration records, and the circuits compiled into native hardware operations.

These resource comparisons have limits. Compound membership, dose grids, and circuit parameters change along with circuit width, so performance differences cannot be attributed to qubit count alone. The pointwise Wilson intervals assume independent shots within an execution and do not capture changes in hardware behavior between runs. A single observed outcome distribution is not enough to predict the success of future hardware batches. Score-evaluation counts also exclude quantum training and execution costs, so they do not establish an end-to-end runtime comparison.

Our analysis extends earlier work that separates feasibility from the quality of valid outputs in trapped-ion optimization \citep{niroula2022} and compares drug-combination models with classical annealing \citep{ramos2026}. Our initial product distribution uses model scores. Guarantees for warm starts based on a continuous relaxation of a binary problem therefore do not automatically apply \citep{egger2021}. The Grover-mixer simulations confirm that the circuit can preserve feasibility \citep{hadfield2019,bartschi2020}, although its improvement in optimum probability over penalty-QAOA depends on the panel. Initialization and training also change in that comparison, so the results do not isolate the effect of the mixer. We construct the initial feasible state directly in simulation; the cost of preparing it on hardware remains to be assessed. Comparisons with classical methods that use the objective score remain necessary \citep{koch2026}.

The drug-response interpretation is also limited by the model. The 647-V panels add a second cell line and show that the CAMA-1 sampling results do not transfer automatically. Overlapping compound panels within a cell line are not independent biological replications, and the original six instances share one optimum. Scaling pairwise synergy by dose fractions extrapolates beyond the measured information used to build the model. The class-weighted dose term is a proxy, not measured toxicity. The resulting score has not been validated as a three-drug response or as an assessment of treatment safety.

\section{Conclusion}
On these drug-response models, producing valid assignments more often than random binary strings did not reliably lead to finding the best assignment. Warm-start QAOA improved some sampling probabilities, but classical greedy search solved every hardware-associated model at the archived score weight, and the results do not establish a computational advantage. Evaluation should report feasibility, optimum recovery, and classical search performance separately.

Future work should test whether longer training changes the results, compare mixers under matched initialization and training conditions, and measure the hardware cost of preparing feasible states. Repeated hardware runs with calibration records would help explain execution shortfalls. Broader cell-line coverage and measured multidrug dose-response data are needed before extending the biological interpretation.

\section*{Data availability statement}
We obtained the drug-response data from the AstraZeneca--Sanger Drug Combination DREAM Challenge dataset \citep{menden2019}, subject to its access conditions. The processed model coefficients, archived hardware outputs, frozen experiment designs, completed simulation and classical baseline results, analysis scripts, and LaTeX sources are available at \url{https://github.com/itsShrizon/warm-start-qaoa-drug-response-reproducibility} (release v1.0.0). The repository includes reproduction commands, software requirements, and file hashes. Supplementary Tables~S1--S11 and Figures~S1--S3 provide the detailed methods and comparisons.


\section*{Author contributions}
TH contributed to methodology, software development, formal analysis, data curation, visualization, and the original draft. RR contributed to conceptualization, methodology, project administration, supervision, funding acquisition, interpretation of results, and manuscript writing and editing. PDB contributed to methodology, validation, interpretation of results, and manuscript review and editing. AAIS contributed to interpretation of the drug-response model and its biological context, and to manuscript review and editing. NH contributed to interpretation of the study, critical review, and editing. PE contributed to high-performance and quantum computing methodology, computational resource support, technical validation, and manuscript review and editing. RS contributed to supervision, research direction, interpretation of the findings, and manuscript review and editing. BWS contributed to conceptualization, methodological guidance, supervision, interpretation of the quantum-computing results, and manuscript review and editing.
% All authors read and approved the final manuscript.



\section*{Funding}
The Pawsey Supercomputing Research Centre supported this work through computational resources allocated under the Setonix-Q Pilot Scheme, project NCMAS-2026-66, ``Closing the Causality Gap: A Computational Framework for Validating Allosteric Drug Mechanisms.'' The allocation provided access to Pawsey's Setonix high-performance computing infrastructure and quantum-enabled computing resources, including quantum processing resources through the Quantum Hub Portal and AWS Braket.


\section*{Conflict of interest}
The authors declare that they conducted the research without any commercial or financial relationships that could be considered a potential conflict of interest.


\section*{Acknowledgments}
The authors have no additional acknowledgments.

\bibliographystyle{Frontiers-Harvard}
\bibliography{references}

\clearpage
\begin{table}[t]
\caption{Hardware results for the original six-task series. Each instance has 200 shots. $q$: qubits; $n$: compounds; $D$: nonzero dose choices; CX: submitted CNOT gates; $|\Fset|$: number of feasible assignments, counted exactly. Parentheses show pointwise 95\% Wilson confidence intervals. The optimum column counts outputs with the globally best feasible score. SA gives the percentage of 100 simulated-annealing runs that reach the optimum within 200 unique score evaluations.}
\label{tab:hardware}
\centering
\small
\setlength{\tabcolsep}{3pt}
\begin{tabular}{rrrrrrlrr}
\toprule
$q$ & $n$ & $D$ & CX & $|\Fset|$ & Valid shots & Valid \% (95\% CI) & Optimum & SA (\%)\\
\midrule
\TableOneRows
\bottomrule
\end{tabular}
\end{table}

\begin{table}[h!]
\caption{Results from the 19 additional hardware tasks. Task creation dates and times are shown in UTC in 2026. $q$: qubits; $n$: compounds; $D$: nonzero dose choices; $N$: shots. Valid-output percentages include pointwise 95\% Wilson confidence intervals. Opt. counts observed globally optimal outputs. Rows for the same instance represent separate completed tasks.}
\label{tab:expanded}
\centering\small
\setlength{\tabcolsep}{5pt}
\begin{tabular}{rrrlrlr}
\toprule
$q$ & $n$ & $D$ & UTC date/time & $N$ & Valid \% (95\% CI) & Opt.\\
\midrule
\RecoveredRows
\bottomrule
\end{tabular}
\end{table}

\begin{figure}[h!]
\centering
\altfigure{figures/figure1.pdf}{Panel A compares feasibility for the original six circuits. Panel B compares individual hardware executions, initial product distributions and ideal circuits through 35 qubits with 4,900 feasible assignments. The last two ideal values are 33.39 and 35.85 percent, compared with hardware values of 6.5 and 3.5 percent. Uniform feasible sampling is always valid.}{\textwidth}
\caption{Feasibility compared with classical controls and noiseless circuit predictions. \textbf{(A)} The original six circuits. \textbf{(B)} The larger circuits, with individual hardware executions offset horizontally and feasible-set sizes shown below the qubit counts. Hardware error bars are pointwise 95\% Wilson confidence intervals. Other series show numerically calculated probabilities; repeated tasks share one ideal-circuit reference. Connecting lines help distinguish the series and do not establish a scaling law.}
\label{fig:feasibility}
\end{figure}

\clearpage
\begin{figure}[t]
\centering
\altfigure{figures/figure2.pdf}{Panel A shows exact optimum probabilities and hardware counts with uncertainty. In the original six-task subset, hardware observes an optimum only at six and eight qubits. Panel B shows calculated 200-draw optimum recovery for the analytic controls. At 15 qubits the rates are 42.7 percent for uniform feasible sampling, 39.0 percent for ideal QAOA, and 25.0 percent for the initial product distribution.}{\textwidth}
\caption{Optimum recovery compared with the sampling controls. \textbf{(A)} Per-shot optimum probability. Hardware error bars are pointwise 95\% Wilson confidence intervals, including an upper bound when no optimum is observed. Markers are offset slightly for visibility. \textbf{(B)} Calculated probability of finding at least one optimum in 200 independent samples from the three analytic distributions, using Equation~\ref{eq:budget}. These calculations assume independent draws.}
\label{fig:optimum}
\end{figure}
\end{document}
